% Generated draft: requires mathematical and source review.
\documentclass[11pt]{article}
\usepackage[margin=25mm]{geometry}
\usepackage{fontspec}
\setmainfont{Latin Modern Roman}
\usepackage{amsmath,amssymb,mathtools,unicode-math}
\setmathfont{Latin Modern Math}
\usepackage{enumitem,needspace,xurl,hyperref}
\setlength{\parindent}{0pt}
\setlength{\parskip}{5pt}
\setlength{\emergencystretch}{3em}
\newcommand{\N}{\mathbb{N}}
\newcommand{\Z}{\mathbb{Z}}
\newcommand{\Q}{\mathbb{Q}}
\newcommand{\R}{\mathbb{R}}
\newcommand{\C}{\mathbb{C}}
\newcommand{\F}{\mathbb{F}}
\newcommand{\A}{\mathbb{A}}
\newcommand{\PP}{\mathbb{P}}
\newcommand{\E}{\mathbb{E}}
\newcommand{\eps}{\varepsilon}
\newcommand{\dd}{\,\mathrm{d}}
\begin{document}
\section*{UnsolvedMath ID: 30002758}
\textbf{Source problem code:} catalogue entry
\subsection{Definitions and mathematical statement}
Let \(x\in\R\) and \(t\geq0\). Consider a real-valued scalar \(\theta(x,t)\) satisfying
\[
\partial_t\theta+u\,\partial_x\theta=0,
\qquad \theta(x,0)=\theta_0(x).
\]
The supplied record compares two velocity laws, corresponding respectively to its models (3) and (4):
\begin{align*}
u_3(x,t)&=\operatorname{PV}\int_{\R}\frac{\theta(y,t)}{y-x}\,\dd y,
\\
u_{4,\gamma}(x,t)&=\int_{\R}|y-x|^{-\gamma}\theta(y,t)\,\dd y,
\qquad 0<\gamma<1.
\end{align*}
Here \(\operatorname{PV}\) denotes the Cauchy principal value at \(y=x\); convergence at infinity must also be ensured. The first law is attributed in the record to C\'ordoba--C\'ordoba--Fontelos, and the second is described there as less singular. The displayed normalizations and signs are retained exactly.

The stated initial-data conditions are symmetry about the origin, nonnegativity, and a maximum at the origin:
\[
\theta_0(-x)=\theta_0(x),\qquad
0\leq\theta_0(x)\leq\theta_0(0)
\quad\text{for every }x\in\R.
\]
The record supplies no additional regularity, decay, integrability, or solution-class hypotheses. The equations are therefore understood only where the indicated integrals and derivatives exist; this does not specify a complete admissible class for the questions below.

The requested first-singularity questions are:
\begin{enumerate}
\item For initial data satisfying the displayed conditions, must the first singularity always be a cusp, or do there exist such initial data producing square-root behavior? The record does not unambiguously assign this question to one or both velocity laws, nor does it define either profile or its asymptotic meaning.
\item For model (4), can singularities occur only where \(\theta'=\partial_x\theta\neq0\), and, for model (3), can they occur only where \(\theta'=0\)? These derivative conditions are reproduced as supplied. Their evaluation time, the relevant profile, and the interpretation of the derivative at a singularity are unspecified.
\item For model (4), construct a solution with a finite first singular time \(T>0\) whose first singularity is an infinite slope of \(\theta\) at \(x=0\), while \(\theta\) remains smooth elsewhere. The record does not specify the limiting meaning of infinite slope or whether smoothness elsewhere refers to a terminal profile, local continuation, or another regularity condition.
\end{enumerate}

These are formulation targets, not asserted conclusions. A precise quantified conjecture requires an admissible initial-data and solution class, a criterion for the first singular time, definitions of cusp and square-root behavior, and a meaning for the derivative-location restrictions. The quantification over \(\gamma\in(0,1)\) and the scope of the initial-data restrictions in the final construction also require clarification. Initial evenness is not asserted to persist under either evolution.
\subsection{Short English statement}
The problem asks how the first singularity forms in two nonlocal transport models: whether its profile must be a cusp or can exhibit square-root behavior, whether its location is constrained by the spatial derivative, and whether the second model admits an isolated first infinite slope at the origin. The source leaves essential regularity and profile definitions unspecified. This is a statement-only draft requiring human mathematical review.
\subsection{Sources}
\begin{itemize}
\item Local UnsolvedMath registry, supplied complete source record: permanent ID \(30002758\), source problem code \texttt{catalogue entry}, title \emph{First-Singularity Profiles in Nonlocal Transport Models}. The record belongs to the Oberwolfach Reports Open Problems collection and gives the proposed year as 2015; no individual proposer is supplied.
\item Supplied report citation: \emph{Spectral Theory and Weyl Functions} (2015), Oberwolfach Reports, DOI 10.4171/owr/2015/1, pp.\ 5--86. Source URL: \url{https://doi.org/10.4171/owr/2015/1}. The supplied background describes these reports as documenting mathematical workshops and their open-problem discussions. No problem-specific page or individual contribution is identified.
\item Supplied literature source, labeled ``original OWR report'': \url{https://doi.org/10.4171/owr/2015/1}.
\item Supplied literature assessment, checked exactly \texttt{2026-08-21}: the record labels the problem open and reports that the cusp, square-root profile, location, and prescribed first-singularity questions remain open in the report, with no primary result located by that cutoff establishing the claims for the stated models. This is supplied dated curation, not an independently verified literature assessment. The source itself requests independent verification of time-sensitive status and cited claims.
\item The registry supplies the statement status \texttt{corrected\_verified} and explains that mathematical content was preserved while notation, typography, or omitted definitions were repaired for publication. This is a supplied editorial label, not independent mathematical verification of the formulation or its status. The velocity attribution to C\'ordoba--C\'ordoba--Fontelos is supplied without a separate bibliographic citation.
\end{itemize}
\end{document}
